\section{Bài dọn dẹp lại cơ chế này}
\label{sec:ch06-cau-noi}

\index{Luong 2015!ba chỗ khác Bahdanau}%
Cơ chế ở mục~\ref{sec:ch06-tron-lai} có vài lựa chọn mà bài đưa ra rồi đi tiếp,
không giải thích và không thử cách khác: chấm điểm bằng \(\vect{s}_{i-1}\) chứ
không phải \(\vect{s}_i\), chấm bằng một mạng một tầng ẩn chứ không phải bằng
\tn{tích vô hướng}{dot product}, và chấm trên mọi vị trí nguồn chứ không phải
một cửa sổ. Một năm
sau, Luong, Pham và Manning thử hết~\autocite{luong2015effective}.

\begin{bridge}{Luong và cộng sự 2015: cùng cơ chế, dọn lại và đo từng lựa chọn}
Bài này đặt tên cho hai họ. \emph{Global} là chấm điểm trên mọi vị trí nguồn,
tức đúng cách bài Bahdanau làm. \emph{Local} là chọn trước một vị trí rồi chỉ
chấm trong một cửa sổ quanh nó, để khỏi trả chi phí trên cả câu.

Mục 3.1 của họ liệt kê ba chỗ khác bài Bahdanau, và tôi chép lại theo lời họ
chứ không tự so:

Thứ nhất, họ dùng trạng thái của tầng LSTM trên cùng ở cả hai phía. Bài
Bahdanau nối trạng thái xuôi và ngược của
\tn{encoder hai chiều}{bidirectional encoder}, còn decoder thì một chiều và
không xếp tầng.

Thứ hai, đường tính khác nhau. Của họ là
\(\vect{h}_t \to \vect{a}_t \to \vect{c}_t \to \tilde{\vect{h}}_t\), tức bước
trước rồi mới chấm điểm bằng trạng thái \emph{mới}. Của Bahdanau là
\(\vect{s}_{i-1} \to \vect{a}_i \to \vect{c}_i \to \vect{s}_i\), rồi còn qua
một tầng deep output và một tầng maxout nữa. Đây đúng là chi tiết mà
mục~\ref{sec:ch06-tron-lai} bảo bạn đọc chậm.

Thứ ba, Bahdanau chỉ thử một hàm chấm điểm, hàm nối hai vector rồi cho qua một
tầng ẩn. Luong thử bốn và đo từng cái. Ba hàm dựa trên nội dung, viết theo ký
hiệu của họ, trong đó \code{h\_t} là trạng thái đích và \code{h\textasciitilde\_s}
là trạng thái nguồn (bài đặt dấu ngang trên chữ cái sau để phân biệt hai bên):

\begin{minted}{text}
dot      h_t^T h~_s
general  h_t^T W_a h~_s
concat   W_a [h_t ; h~_s]
\end{minted}

cộng một hàm chỉ nhìn vị trí. Kết quả ở bảng 4 của họ, trên WMT'14 Anh-Đức:
\code{dot} hợp với global, \code{general} hợp với local. Còn \code{concat}, tức
họ hàng gần nhất của hàm Bahdanau dùng, là hàm họ không làm cho chạy tốt được.
Nguyên văn: \enquote{our implementation of concat does not yield good
performances and more analysis should be done to understand the reason}.

Đó là một câu nên đọc theo cả hai chiều. Nó không nói hàm của Bahdanau tệ, vì
Bahdanau dùng nó và đạt kết quả tốt. Nó nói rằng lựa chọn hàm chấm điểm là một
lựa chọn thật, đủ sức làm hỏng một mô hình, và tới 2015 vẫn chưa ai hiểu vì
sao. Bài kế tiếp trong sách này chọn tích vô hướng, và
chương~\ref{ch:transformer} có nguyên một dẫn xuất về chỗ tại sao tích vô hướng
cần chia cho \(\sqrt{\dk}\) mới dùng được.

Bài cũng đo \tn{đồng chỉnh}{alignment} bằng alignment error rate trên dữ liệu
người làm, và
đó là nguồn của nhận xét ở mục~\ref{sec:ch06-mo-ma-tran}:
\tn{điểm đồng chỉnh}{alignment score} và điểm dịch không đi cùng nhau. Kết quả
tổng: attention đem lại tới 5.0 BLEU so
với hệ không attention đã có sẵn dropout và phép đảo câu nguồn.

EMNLP 2015, trang 1412-1421~\autocite{luong2015effective}. Tôi đọc trang 1412
tới 1419; hai trang cuối là bảng câu dịch mẫu và chương này không dùng tới.
\end{bridge}

\index{Luong 2015!local attention và chi phí}%
Nửa \emph{local} của bài Luong đáng chú ý vì một lý do không liên quan gì tới
chất lượng dịch. Nó là phản ứng đầu tiên tôi biết trước cái giá mà bài Bahdanau
đã tự ghi nhận: chấm điểm mọi vị trí nguồn cho mọi từ đích. Luong giải bằng
cách không chấm hết, chỉ chấm trong một cửa sổ.

Cách ấy không phải cách còn được dùng hôm nay. Cái còn lại, mười năm sau, là
giữ nguyên phép chấm điểm đầy đủ và viết lại nó cho hợp với bộ nhớ của phần
cứng, và chương~\ref{ch:chi-phi-bac-hai} là chương kể chuyện đó. Nhưng năm 2015 thì câu
hỏi đã được đặt ra rồi, và mục sau là chỗ bài Bahdanau tự đặt nó.
